% load package if document contains at least one citation
\usepackage[maxbibnames=99,sorting=none,style=numeric,backend=biber]{biblatex}
% enable option to print bibliography as a numbered subsection
\defbibheading{bibsubsection}[\bibname]{%
  \subsection{#1}%
}
\addbibresource{Private/PHBern/References.bib}

\begin{document}

\lhead[ \leftmark ]{PHBern, Grundlagen der Didaktik}
\rhead[Informatik]{Cyril Wendl, \today, Bern}


\section*{Themenvertiefung \enquote{Lernprozesse verstehen}}

Im Folgenden möchte ich auf Aspekte aus den Texten \parencite{alma99116838938605511,alma99116996690305511,alma99117228235505511} eingehen, welche ich für meinen eigenen Unterricht mitnehme. Jede der folgenden Boxen fasst jeweils die relevante zugrundeliegende Theorie kurz und stichwortartig zusammen und schliesst danach mit persönlichen Schlussfolgerungen ab.

\begin{myExBox}[title=Intermettierendes Verstärken von gewünschtem Lern-Verhalten]
\subsubsection*{Theoretische Grundlage}
Im Artikel von \textcite{alma99117228235505511} wird auf unterschiedliche Lerntheorien eingegangen, vom Behaviourismus bis zu kognitiven Richtungen. Im Rahmen unterschiedlicher behaviouristischer Strömungen erwähnen die Autoren unter anderem 
das \textbf{operante Lernen nach Skinner}: Skinner hat aufgezeigt, dass Organismen nicht nur reaktiv sind (im Gegensatz zur klassischen Theorie des Konditionierens nach Pawlow), sondern auch von innen heraus operieren --- beispielsweise mit dem \textit{Skinner-Box}-Experiment, bei dem ein Stück Futter nur dann in einen Tierkäfig fällt, wenn das Tier einen bestimmten Hebel drückt. Das Tier lernt also durch \enquote{\textit{trial and error}}, was funktioniert und was nicht, woraus Skinner die Begriffe des \textit{Lernens am Erfolg (law of effect)} sowie der Verbindung von Stimulus und Reaktion (S-R-Verbindung) formulierte.
An der Theorie wurden Erweiterungen und Differenzierungen vorgenommen, etwa um zu erforschen, welche Mechanismen zu einer Verstärkung oder Löschung der bereits erlernten Reaktionen bzw. Lernstrategien führen können. Insbesondere wurde herausgefunden, dass intermettierendes, also unregelmässiges Verstärken (z.B. Futterkorn wird nur jedes dritte Mal gegeben bzw. Belohnung kommt nicht immer) zu einer noch stärkeren Resilienz gegen Löschung führen kann (dieser Effekt wird übrigens auch von sozialen Netzwerken ausgenutzt, um Nutzerinnen und Nutzer stärker an sich zu binden).

\subsubsection*{Anwendungen auf eigenen Unterricht}
Im Unterricht begegnet mir häufig die Herausforderung, dass SuS bereits nach kurzer Zeit neue Lerninhalte wieder vergessen haben. Hierzu würde mich interessieren, inwiefern das \textit{intermettierende Verstärken und Wiederholen} gewissern Lerninhalte, bzw. Lernstrategien, welches Skinner im Rahmen seiner Verstärkungspläne als besonders resistent gegen die Löschung erachtet, nützlich sein könnte: Beispielsweise könnte ich die SuS intermittierend in Ihren Bemühungen bekräftigen, Ihre Hausaufgaben zu machen (Verhaltensmuster fördern, positives Feedback), oder gewisse Inhalte (z.B. Grundkenntnisse des Programmierens) intermittierend repetieren und eventuell mit einem Anreiz (spontane Mini-Noten o.ä.) verbinden.
\end{myExBox}

\begin{myExBox}[title=Sinnvolles Lernen mit \textit{Advance Organizer}]
\subsubsection*{Theoretische Grundlage}
Im Rahmen der kognitiven Lern-Theorien erwähnen  \textcite{alma99117228235505511} die Wichtigkeit, Lernen nicht nur als mechanisches (Auswendig-)Lernen zu betrachten, sondern den Aufbau von kognitiven Strukturen im Auge zu behalten. Die kognitive Psychologie beschäftigt sich insbesondere mit Lerntheorien, bei denen das Verhalten der Lernenden Hinweise auf Vorgänge im Gedächtnis des Menschen gibt. Die \textit{Innensteuerung} von Lernprozessen wird also betont. Beispiele von Forschungs-Erkenntnissen der kognitiven Psychologie sind: \begin{enumerate*}
\item Das Hirn ist kein \enquote{Ablagemechanismus}, sondern arbeitet in komplexen Netzen, die \textbf{Vernetzung} mit dem Vorwissen ist also zentral.
\item \textbf{Der emotionale Kontext} spielt eine zentrale Rolle für die Lernfähigkeit (wie, wann, unter welchen Umständen hat ein Lern-Ereignis stattgefunden?)
\item Das Hirn arbeitet in neuronalen Netzen.
\item Gedächtnis und Bedeutung durch den / die Empfänger(in) hängen eng zusammen
\end{enumerate*}, um nur einige Beispiele zu nennen. 

\subsubsection*{Anwendungen auf eigenen Unterricht} Insbesondere die Wichtigkeit und den Sinn des Gelernten, also die Sinnesstrukturen zu vermitteln, scheint mir zentral, um tiefes Lernen zu ermöglichen (also kein kurzfristiges \enquote{Bulimie-Auswendiglernen}). \textcite{alma99117228235505511} schlagen diesbezüglich vor, einen Advance Organizer zu verwenden, welcher folgenden Aspekten dient:
\begin{enumerate*}
\item Mobilisierung der Vorkenntnisse der Lernenden
\item Sinnvolle Verknüpfungen zwischen schon vorhandenem und neuem Wissen herstellen
\item Verstehen anbahnen.
\end{enumerate*}
Einen Advance Organizer verwende ich zwar regelmässig, aber nicht systematisch. Ich werde in Zukunft versuchen, diesen noch häufiger im Unterricht zu verwenden, um die Sinnhaftigkeit der vermittelten Inhalte und deren Einbettung im grösseren Kontext aufzuzeigen.
\end{myExBox}

\begin{myExBox}[title=Emotionen verstehen und positiv beeinflussen]
\subsubsection*{Theoretische Grundlage}
Zwei zentrale Aussagen des Texts von \textcite{alma99116838938605511} sind, dass Emotionen in Lehr- und Lernprozessen omnipräsent sind und dass sie beide Prozesse massgeblich beeinflussen. Es liegt also nahe, dass eine Auseinandersetzung mit und ein Verständnis von Emotions-Theorien und -Prozessen zentral für den Wissensvermittlungsprozess sind. Mit Emotionen sind \enquote{episodische psychische Zustände von Personen} gemeint, wie z.B. Freude, Traurigkeit, Angst oder Ärger, welche meist in die Kategorien angenehme und unangenehme Gefühle eingeteilt werden können. Die \textit{appraisal theory} ist die dominierende Theorie der Emotionsentstehung in der heutigen Psychologie und will die Entstehung sämtlicher Emotionen des Menschen erklären, indem angenommen wird, dass Einschätzungen (\textit{appraisals}) von Sachverhalten Emotionen hervorrufen, und dass Menschen zielgesteuerte Wesen sind. Der Erfolg oder Misserfolg im Erreichen von gewissen Zielen kann daher positive oder negative Emotionen wecken. Allerdings gibt es auch nicht-selbstbezogene, sogenannte \enquote{soziale} Emotionen wie beispielsweise die Mit-Freude oder das Mit-Leid am (Miss-)Erfolg einer anderen Person --- in diesem Fall spricht man von \textit{sozialen Emotionen}. Nicht zuletzt haben Emotionen diverse Funktionen, von denen drei hervorgehoben werden: \begin{enumerate*}
\item Die Aufmerksamkeit fokussieren
\item Informieren (z.B. \enquote{was ist gerade passiert?})
\item Motivieren.
\end{enumerate*}

\subsubsection*{Anwendungen auf eigenen Unterricht}
Mir scheint wichtig, die \textit{Funktionen von Emotionen} im Unterricht im Auge zu behalten, also beisipelsweise die positiven, adaptiven Funktionen von Emotionen auf das Lernen besser zu verstehen und ggf. zu nutzen. Was kann beispielsweise getan werden, wenn SuS nach einer schwierigen Prüfung in der Folgelektion nicht fähig sind, neue Inhalte aufzunehmen? Kann ich negative Emotionen allenfalls transformieren, indem ich die Prüfung nachbespreche, den SuS Zeit zum Reflektieren gebe oder eine Möglichkeit biete, Rückfragen zur Prüfung zu stellen? Wie kann ich danach mittels neuer Emotionen den Fokus darauf lenken, die SuS für ein neues Thema zu motivieren? Um dies zu erreichen, könnte ich allenfalls, je nach SuS, individuelle Motivationsaspekte aus der hedonistischen Theorie der Motivation (z.B. \enquote{das kannst Du damit erreichen / verhindern}) oder aus der Theorie der emotionalen Handlungsimpulse (nicht-hedonistische Handlungswünsche, z.B. \enquote{damit kannst Du Dich nützlich machen / anderen helfen}) mobilisieren.\end{myExBox}

\begin{myExBox}[title=Emotionen zur Fokussierung von Aufmerksamkeit]
\subsubsection*{Theoretische Grundlage}
Wie bereits erwähnt gehen \textcite{alma99116838938605511} auf die Funktion von Emotionen ein: Insbesondere wird in Kapitel 6.3.1 die Emotions-Funktion ausgeführt, die Aufmerksamkeit zu fokussieren, d.h. die Aufmerksamkeit auf das Ereignis zu lenken, welches die Emotionen ausgelöst hat. Insbeondere bei negativen und überraschenden Ereignissen ist der Nutzen dieser Fokussierung offensichtlich. Umfassende Forschung zu Thema Überraschung hat gezeigt, dass Überraschung nicht nur die Aufmerksamkeit auf ein Thema lenken kann, sondern auch tiefere Verarbeitungsprozesse anstösst, insbesondere betreffend der Ursache des überraschenden Ereignisses.

\subsubsection*{Anwendungen auf eigenen Unterricht}
Gerne würde ich im Unterricht häufiger mit dem Überraschungs-Effekt experimentieren, um die Aufmerksamkeit auf ein (neues) Thema zu fokussieren, die Auseinandersetzungstiefe zu steigern und somit für eine allgemein packendere Lernatmosphäre zu sorgen. Gerade in der Informatik gibt es viele überraschende Beispiele und ''Zauber-Experimente'', die zu diesem Effekt führen könnten.
\end{myExBox}

\begin{myExBox}[title=Interferenzen vermeiden]
\subsubsection*{Theoretische Grundlage}
\textcite{alma99116996690305511} geht auf die Psychologie des Lernens und des Lehrens ein, indem er zuerst beim Gehirn ansetzt und dessen unterschiedliche Bestandteile sowie deren Interaktionen von neurologischer Seite her zusammenfasst. Das Gedächtnis wird als \enquote{kognitives System, das Informationen aufnimmt, enkodiert, modifiziert und wieder abruft} beschrieben. Das Gedächtnis kann aufgeteilt werden in ein \begin{enumerate*}
\item Ultrakurzzeitgedächtnis (UKZG, sensorische Register für äussere Reize)
\item Kurzzeitgedächtnis (KZG)
\item Langzeitgedächtnis (LZG)
\end{enumerate*}, wobei die drei \enquote{Gedächtnis-Typen} alle unterschiedliche Kapazitäten und Informations-Zerfallsraten haben. Im Folgenden wird auf verschiedene Inferenztypen eingegangen, die für das Abspeichern neuer Inhalte hinderlich sein können.
\subsubsection*{Anwendungen auf eigenen Unterricht}
Im eigenen Unterricht möchte ich mich vermehrt darauf achten, ob Interferenzen inhaltlicher sowie anderer Arten die Lerninhalts-Vermittlung erschweren. Dabei kommen mir insbesondere \textit{Ähnlichkeitshemmungen} in den Sinn, also die Tatsache, dass Inhalte schwerer reproduziert werden, wenn bereits ein ähnlicher Inhalt gespeichert ist. Die passiert etwa beim Thema des Gleichzeichens (\enquote{=}), welches in der Mathematik eine andere Bedeutung hat als in der Informatik (Gleichheit vs. Zuweisung). Solche Hemmungen möchte ich wenn immer möglich identifizieren und beseitigen, um Klarheit zu schaffen.
\end{myExBox}


\begin{myExBox}[title=\textit{Dual-Code}-Theorie \& Multi-Modalität im Unterricht]
\subsubsection*{Theoretische Grundlage}
\textcite{alma99116838938605511} weisen auf die \textit{Dual-Code}-Theorie hin, welche auf neuro-psychologischen Studien basiert, die anhand empirischer Evidenzen zur Einsicht kamen, dass bei einer doppelten Kodierung von Inhalten (z.B. Bild \textit{und} Text) die Inhalts-Behaltensrate besonders hoch ist.
\subsubsection*{Anwendungen auf eigenen Unterricht}
In meinem Unterricht achte ich mich bereits jetzt darauf, neue Inhalte nicht nur zu erklären, sondern, wenn immer möglich, auch mittels Bildern, Vergleichen, Metaphern, Videos, schematischen Darstellungen usw. in unterschiedlicher Weise darzustellen. Dieser Aspekt scheint mir insbesondere bei stark theoretischen Inhalten sehr wichtig, daher werde ich in Zukunft noch konsequenter auf die \textit{Dual-Code}-Theorie setzen und mich damit auseinandersetzen, welche anderen Medien in meinem Unterricht noch förderlich sein könnten.
\end{myExBox}

\end{document}